\begin{center}
\begin{tcolorbox}[colback=gray!5,colframe=gray!40]
\textit{This chapter has been published as part of in F. Adesanya et al., ``Systems of twinned systems: A systematic literature review''~\cite{adesanya2025systems}.}
\end{tcolorbox}
\end{center}

SoTS emerge from the intersection of SoSn and DTs. SoS research focuses on coordinating autonomous, heterogeneous systems to achieve collective goals~\cite{maier1998architecting}, while DTs provide continuously synchronized digital representations of physical assets capable of real-time monitoring, prediction, and control~\cite{kritzinger2018digital, tao2018digital}. Integrating these two ideas enables new forms of coordination: DTs offer SoS stronger situational awareness and predictive capability, while SoS provide DTs a scalable, multi-system context in which collective behaviour can emerge.


To clarify how this convergence is currently understood and implemented, we conducted a systematic literature review following the study design described in \secref{sec:literature-review-study-design}. The review also introduces a SoTS classification framework (\secref{sec:literature-review-classification}) that structures the analysis of architectural patterns and constituent relationships. This chapter synthesizes the findings most relevant to this thesis. 
\secref{subsec:literature-review-results-motivations} examines why SoTS are formed, including the motivations, intents, and domains in which DTs and SoS are combined. \secref{subsec:literature-review-results-architecture} analyzes how SoTS are structured by surveying the dominant coordination patterns and architectural types. \secref{subsec:literature-review-results-sos-characteristics} evaluates which SoS characteristics are realized in practice, with particular attention to dynamic properties such as emergence, reconfiguration, and evolution. \secref{subsec:literature-review-results-reliability} then considers how non-functional properties, specifically reliability, are addressed. Finally, \secref{subsec:literature-review-results-standards} and \secref{subsec:literature-review-results-trl} discuss the maturity of existing approaches, including their adoption of standards and their reported technological readiness levels (TRLs), highlighting the extent to which SoTS have been evaluated beyond controlled or simulated settings.
